%% kapa: exam=1 round=37 year=2026 period=1 form=A subject=real_estate subject_ko=부동산학원론 seq=12 number=92
%% source: Q-Net 감정평가사 기출문제 (cst00309) artlSeq=5258924 file=raw/1008/5258924/2026년 제37회 감정평가사 1차 1교시.pdf
%% conversion: kordoc 4.24.51 -> tools/gen1.py (auto); review=proofread-sonnet; flags=none
\begin{question}{92}
디파스퀠리-위튼(DiPasquale \& Wheaton)의 사분면 모형에 관한 설명으로 옳은 것은? (단, 주어진 조건에 한함)
\begin{choices}
  \item 임대료, 자산가격, 신규건설비용, 공간재고 등 4개 내생변수의 장기균형이 결정된다.
  \item 지역경제변수인 소득, 인구 또는 고용이 증가하면 부동산공간시장에서 임대료는 단기적으로 하락한다.
  \item 부동산자산시장에서 자산가격은 자본투자에 대한 요구수익률을 의미하는 내생변수이다.
  \item 건물의 신규건설량과 최초 공간재고가 일치할 때 장기균형에 도달한다.
  \item 건축원자재 가격이 하락하면 새로운 장기균형에서 균형공간재고는 증가한다.
\end{choices}
\end{question}
\answer{5}
